% Lab part of the module.
%
% Module MC-3 of Section 6.1 of the syllabus: "Spectral features for motion
% data". Type T+L, 90 minutes: about 45 minutes of theory and 45 minutes of
% lab. The lab needs no board. It runs on the recordings of a group, or on
% the simulated fallback dataset of its folder.
%
% The parts follow the README of the lab folder. The frame of the expected
% output shows the form of the output and the values that check the method.
% No value here is a measurement of a board.

% --------------------------------------------------------------- 1. the goal
\begin{frame}{Lab: goal and plan}
  \small
  \begin{itemize}
    \item \textbf{Goal.} You compute the features of a motion signal, train
          a small classifier on them, and compare five inputs for the same
          task.
    \item \textbf{Deliverable.} The table of section 8 of the notebook,
          copied into \code{report.md}, and the answers of its three questions.
  \end{itemize}
  \vskip 0.5em
  \footnotesize
  \setlength{\tabcolsep}{4pt}
  \begin{tabular}{@{}L{1.0cm}L{7.2cm}L{1.5cm}@{}}
    \toprule
    Part & Content & Time \\
    \midrule
    A & The data, the windows, and the three time-domain features & 12 min \\
    B & The FFT and the spectral power of one axis & 13 min \\
    C & The classifier, and its number of parameters & 8 min \\
    D & Five inputs for the same task, and the FFT length & 12 min \\
    \bottomrule
  \end{tabular}
  \vskip 0.5em
\end{frame}

% ------------------------------------------------------------ 2. the files
\begin{frame}{Lab: files and data}
  \small
  \setlength{\tabcolsep}{4pt}
  \begin{tabular}{@{}L{3.9cm}L{5.3cm}@{}}
    \toprule
    File & Content \\
    \midrule
    \code{spectral\_features.ipynb} & The notebook. It has four tasks. \\
    \code{host/motion\_sim.py} & The simulated signals. \\
    \code{host/make\_dataset.py} & Makes the fallback dataset. \\
    \code{report.md} & The report to hand in. \\
    \code{solutions/} & The notebook with all tasks complete and with its output, and an example report. \\
    \bottomrule
  \end{tabular}
  \vskip 0.5em
  \small
  \begin{itemize}
    \item The notebook reads the folder \code{data/} when it has recordings of
          the format of the course. Otherwise it makes the simulated
          dataset in \code{fallback\_data/}.
    \item The last session is the test session. A model that separates the
          sessions has learned the person, not the motion.
  \end{itemize}
\end{frame}

% ---------------------------------------------------------------- 3. part A
\begin{frame}[fragile]{Part A: the data and the time features (12 min)}
  \small
  \begin{enumerate}
    \item Start Jupyter from the folder of the lab:
\begin{lstlisting}[style=shell]
../../../.venv/bin/jupyter lab spectral_features.ipynb
\end{lstlisting}
    \item Run section 1. Write the number of training windows and of test
          windows in the report.
    \item \textbf{Task 1.} Write \code{time\_features(window)}: the RMS, the
          skewness, and the kurtosis of one axis, after the removal of the
          mean.
    \item Run the cell after the task.
  \end{enumerate}
  \vskip 0.5em
  \small You see the check of the task, and the three values for the first
  window:
  \begin{center}
  \begin{tabular}{@{}lll@{}}
    \toprule
    RMS & Skewness & Kurtosis \\
    \midrule
    0.0038 & $-0.3266$ & $-0.0771$ \\
    \bottomrule
  \end{tabular}
  \end{center}
  \small The first window is an idle window, so its values are small. That
  is the answer of the check, not a target for your own window.
  \source{Values of one run of the solution notebook on the fallback
    dataset, an x86 processor with AVX2.}
\end{frame}

% ---------------------------------------------------------------- 4. part B
\begin{frame}{Part B: the FFT and the spectral power (13 min)}
  \small
  \begin{enumerate}
    \item Run section 3. The cell draws the spectral power of one window of
          each class on the z axis. The z axis of \code{terrestrial} is almost
          still, because the motion of a truck sits in the horizontal plane.
          Write in the report which frequency holds the power of the class
          \code{lift}, and which one holds the power of \code{maritime}.
    \item \textbf{Task 2.} Write \code{axis\_features(window, nfft)}: the three
          time-domain features, then the power of bins 1 to \code{nfft // 2}.
    \item Run the cell after the task, then section 4.
  \end{enumerate}
  \vskip 0.5em
  \small You see the number of values of one axis, and it must be
  $3 + 128 / 2 = 67$. Then the feature matrix has 201 columns and the
  notebook prints the mean and the standard deviation of the first nine
  features.
  \vskip 0.3em
\small The first three values of the feature vector are the RMS, the
skewness, and the kurtosis of the axis `x`. The next six are power values.
\end{frame}

% ---------------------------------------------------------------- 5. part C
\begin{frame}{Part C: the classifier (8 min)}
  \small
  \begin{enumerate}
    \item \textbf{Task 3.} Write \code{make\_model(n\_inputs)}: a layer of 20 units
          with ReLU, a layer of 10 units with ReLU, and 4 outputs with
          softmax.
    \item Run the cell after the task. The check calculates the expected
          number of parameters itself.
    \item Run the next cell to train it.
  \end{enumerate}
  \vskip 0.5em
  \begin{center}
  \begin{tabular}{@{}l@{}}
    \toprule
    input values 201, parameters 4294 (expected 4294), outputs 4 \\
    the model on the spectral features: 4294 parameters, \\
    test accuracy 98.8 percent \\
    \bottomrule
  \end{tabular}
  \end{center}
  \small The names of the layers and of the input matter. Without them, the
  file of the model changes size with each model that the notebook builds.
  \source{Values of one run of the solution notebook on the fallback
    dataset, an x86 processor with AVX2.}
\end{frame}

% ---------------------------------------------------------------- 6. part D
\begin{frame}{Part D: five inputs for the same task (12 min)}
  \small
  \begin{enumerate}
    \item Run section 6. It trains the same network on the raw window, on
          the time-domain features, and on the spectral features.
    \item Write the three lines of output in the report.
    \item \textbf{Task 4.} Before you run section 7, complete the table with
          the number of values of one axis for the FFT lengths 32, 64, and
          128. Then write the length that you expect to win, and why.
    \item Run section 7 and section 8, and copy the table into the report.
  \end{enumerate}
  \vskip 0.4em
  \small The notebook prints the same table as the figure of the lecture:
  the raw window reaches 58.5 percent, the nine time-domain values reach
  100.0 percent, and the spectral values reach 98.8 percent with an FFT
  length of 128.
  \source{Values of one run of the solution notebook on the fallback
    dataset, an x86 processor with AVX2. The signals are simulated.}
\end{frame}

% ------------------------------------------------------------ 7. the check
\begin{frame}{Check criterion and Decision Log}
  \small
  \begin{itemize}
    \item The notebook prints \code{complete} for tasks 1, 2, 3, and 4.
    \item The table of the report has five rows, each with the number of
          values, the number of parameters, and a measured test accuracy.
    \item The three questions of the report have answers with numbers from
          your own run.
    \item The Decision Log names one decision and one number.
  \end{itemize}
  \vskip 0.5em
  \keypoint{Question of the Decision Log: a device must classify the motion
    of a pallet for a year on a battery. Which feature set do you put in
    the firmware, and which number of your table decides it?}
\end{frame}

% ------------------------------------------------------- 8. common problems
\begin{frame}{Common problems}
  \small
  \begin{tabular}{@{}L{4.4cm}L{5.0cm}@{}}
    \toprule
    Problem & Fix \\
    \midrule
    The notebook says that it uses the fallback dataset & Your group has no
      recording in \code{data/}. The signals are simulated. \\
    \code{Task 2} says that the answer needs 67 values & Your answer has a
      different number. Return one value for each bin. \\
    \code{Task 3} says that the model raises an error & A layer is missing, or a
      layer is added with a tensor instead of with its input. \\
    The accuracies change between two runs & Write the seed, and run the
      same cells again. \\
    \bottomrule
  \end{tabular}
  \vskip 0.5em
  \small The README of the lab folder has the full list of the checks of the
  notebook and the run time.
\end{frame}